% !TeX root = ../机器学习笔记.tex
% 由主文档通过 \input 载入；请编译主文档。
\setcounter{chapter}{10} % 保留原书章号。
\chapter{决策树}
\label{chap:ml-trees}

决策树（decision tree）通过一连串问题把输入空间划分成区域，再为每个区域给出预测。它能表示非线性关系，但训练通常采用贪心搜索，表示能力和找到最优结构的能力必须分开理解。

\section{递归划分与树结构}

内部节点按条件把样本分入子节点，叶节点输出类别分布或数值。对连续特征，常用条件 $x_j\leq t$；沿根到叶的路径给出一组约束，因此叶对应一个输入区域 $R_m$。回归树预测可写为 $\sum_m c_m\mathbf1\{x\in R_m\}$，是分段常数函数。

训练从根节点全部样本开始，选择划分，对子节点递归，直到满足停止条件。预测只沿一条路径到叶。训练中的最佳当前划分不一定通向全局最佳整棵树，完全搜索结构通常远比一次局部比较复杂。

\section{分类树的划分准则}

节点中类别比例为 $p_c$，熵为 $H=-\sum_cp_c\log p_c$，Gini 不纯度为 $G=1-\sum_cp_c^2$。若候选划分产生子节点 $v$，其不纯度下降是
\[
 \Delta I=I(\text{父节点})-\sum_v\frac{n_v}{n_{\rm parent}}I(v).
\]
按子节点样本比例加权，才能反映所有样本的平均变化。信息增益使用熵，信息增益率进一步除以子节点比例分布的熵；分母为零的无效划分必须排除。

多取值特征可能通过产生许多很小的纯节点获得很高训练增益，故纯度不是唯一目标。最小叶样本数、验证评价和剪枝限制这种碎片化；不应把不同树算法的划分准则和后处理步骤当作相同流程。

\section{回归树与连续特征切分}

平方损失下，一个叶区域 $R$ 的最优常数为 $\bar y_R$，由对 $\sum_{i\in R}(y_i-c)^2$ 求导得到。候选二分的得分是两个子节点的最小残差平方和之和；选其最小者。绝对损失下叶常数改为中位数，因此输出规则依赖损失。

对一个数值特征先排序，只需考虑相邻不同值之间的阈值；相同值不能被该阈值分开。通过累计样本数、响应和与平方和，能快速计算每个候选的残差平方和，避免为每个阈值重新扫描全部样本。

\section{停止条件与剪枝}

预剪枝用最大深度、最小叶样本数或最小增益限制生长。后剪枝先得到较大的树，再比较删去子树的结果。代价复杂度准则为 $\widehat R(T)+\alpha|L(T)|$，其中 $L(T)$ 为叶集合，$\alpha\geq0$ 控制结构复杂度；$\alpha$ 由验证数据选择。

树常有较高方差：早期切分因数据扰动而改变后，后续节点接收的样本也随之改变。剪枝降低这种敏感性，却可能丢失有用细节。训练时每次切分都降低损失，不代表无限加深后的测试风险会下降。

\section{缺失特征、解释与局限}

缺失值可先填补并增加指示变量，或由具体树算法学习缺失方向；两种处理对应不同假设。树的路径便于解释个体预测，但“容易写出规则”不意味着规则稳定或有因果意义。

分裂重要性累计某特征带来的训练不纯度下降，可能偏好候选切分较多的特征。置换重要性衡量在评价数据上打乱某特征导致的指标变化，相关特征会影响其解释，而且置换可能产生不自然的组合。回归树叶常数通常不能延续训练区域以外的趋势，故外推能力有限。

\section{例题、练习与学习检查}

\begin{example}[一个回归树切分]
输入为 $0,1,2,3$，响应为 $0,0,2,2$，只允许一个阈值。
\end{example}
\begin{solve}
未切分时叶均值为一，残差平方和为四。取阈值 $1.5$ 后，左、右叶均值分别为零和二，总残差为零。阈值 $0.5$ 或 $2.5$ 的总残差均为 $8/3$，因此本次贪心选择为 $1.5$。
\end{solve}

习题：计算两个各含一半样本、类别分别纯净的子节点相对于均衡父节点的熵增益；推导叶均值公式；在上例增加一个异常响应，比较切点和叶值变化；解释训练纯度与可泛化解释之间的差别。





